\documentclass[11pt,a4paper]{article}

\usepackage[utf8]{inputenc}
\usepackage[T1]{fontenc}
\usepackage{lmodern}
\usepackage[margin=2cm]{geometry}   % tighten toward 1.5cm if you run over 3 pages
\usepackage{amsmath,amssymb}
\usepackage{mathpartir}
\usepackage{listings}
\usepackage{xcolor}
\usepackage[hidelinks]{hyperref}

\lstset{
  basicstyle=\ttfamily\small,
  keywordstyle=\bfseries,
  morekeywords={Func,Print,Read,Lambda,Var,While,If,Else,Ret,Nat,Int,Real,Bool,True,False},
  columns=fullflexible,
  frame=single,
}

% --- Front matter details -------------------------------------------------
\newcommand{\langname}{Numa}                    
\newcommand{\studentid}{2763839}
\newcommand{\giturl}{https://github.com/DavidTilman/plpdi-mini-project}
\newcommand{\githash}{TODO}                     % TODO: commit hash to be assessed

% --- Semantics notation (borrow rules from numa; these macros match) ------
\newcommand{\kw}[1]{\mathtt{#1}}                % keyword / concrete syntax
\newcommand{\eval}{\Downarrow}                  % big-step evaluation
\newcommand{\config}[1]{\langle #1 \rangle}     % configuration <e, rho, sigma>
\newcommand{\clo}[3]{\mathsf{clo}(#1,\, #2,\, #3)} % closure
\newcommand{\rn}[1]{\TirNameStyle{#1}}          % rule name

% --- Type Rules ---%
\newcommand{\TENAT}{
	\inferrule*[right=TE-NAT]
	{ }
	{\Gamma\vdash n : \kw{Nat}} 
}

\newcommand{\TEREAL}{
	\inferrule*[right=TE-REAL]
	{ }
	{\Gamma\vdash n : \kw{Real}} 
}

\newcommand{\STREF}{
	\inferrule*[right=S-REF]{ }{t \sqsubseteq t}
}

\newcommand{\STTRANS}{
	\inferrule*[right=S-TRANS]{t \sqsubseteq t' \\ t' \sqsubseteq t''}{t \sqsubseteq t''}
}

\newcommand{\STNATINT}{
	\inferrule*[right=S-NatInt]{ }{\kw{Nat} \sqsubseteq \kw{Int}}
}

\newcommand{\STINTREAL}{
	\inferrule*[right=S-IntReal]{ }{\kw{Int}\sqsubseteq \kw{Real}}
}

\newcommand{\TETRUE}{
	\inferrule*[right=TE-TRUE]
	{ }
	{\Gamma\vdash \kw{True} : \kw{Bool}}
}

\newcommand{\TEFALSE}{
	\inferrule*[right=TE-FALSE]
	{ }
	{\Gamma\vdash \kw{False} : \kw{Bool}}
}

\newcommand{\TEVAR}{
	\inferrule*[right=TE-VAR]
	{x : \tau \in \Gamma}
	{\Gamma\vdash x : \tau}
}

\newcommand{\TELAMBDA}{
	\inferrule*[right=TE-LAMBDA]
	{\Gamma, x_1:\tau_1, \ldots, x_n:\tau_n \vdash e : \tau' \\ \tau' \sqsubseteq \tau}
	{\Gamma \vdash \kw{Lambda}\,(x_1:\tau_1, \ldots, x_n:\tau_n):\tau \Rightarrow e : (\tau_1, \ldots, \tau_n) \to \tau}
}

\newcommand{\TENEGATE}{
	\inferrule*[right=TE-NEGATE]
	{\Gamma\vdash e : \tau \\ \mathsf{neg}(\tau) = \tau'}
	{\Gamma\vdash -e : \tau'} 
}

\newcommand{\TEOP}{
	\inferrule*[right=TE-OP]
	{\Gamma\vdash e_1 : \tau_1 \\ \Gamma\vdash e_2 : \tau_2 \\ \mathsf{sig(\oplus, \tau_1, \tau_2) = \tau}}
	{\Gamma\vdash e_1 \oplus e_2 : \tau}
}

\newcommand{\TEAPP}{
	\inferrule*[right=TE-APP]
	{ \Gamma \vdash e : (\tau_1, \ldots, \tau_n) \to \tau \\ \Gamma \vdash e_i : \tau'_i \\ \tau'_i \sqsubseteq \tau_i \quad (1 \le i \le n) }
	{ \Gamma \vdash e(e_1, \ldots, e_n) : \tau }
}

\newcommand{\TECOND}{
	\inferrule*[right=TE-COND]
	{ \Gamma \vdash e : \kw{Bool} \\ \Gamma \vdash e_1: \tau_1 \\ \Gamma \vdash e_2: \tau_2 \\ \tau = \mathsf{widest}(\tau_1, \tau_2)}
	{\Gamma\vdash (e\ ?\ e_1\ :\ e_2) : \tau}
}

\newcommand{\TSDECL} {
	\inferrule*[right=TS-DECL]
	{ \Gamma \vdash e : \tau' \\ \tau' \sqsubseteq \tau}
	{ \Gamma \vdash_{\tau_r} \kw{Var}\ x:\tau = e \dashv \Gamma, x:\tau}
}

\newcommand{\TSREAD} {
	\inferrule*[right=TS-READ]
	{ \text{Read input is of type }\tau}
	{ \Gamma \vdash_{\tau_r} \kw{Var}\ x : \tau= \kw{Read} \dashv \Gamma,x:\tau}
}

\newcommand{\TSASSIGN} {
	\inferrule*[right=TS-ASSIGN]
	{ x:\tau \in \Gamma \\ \Gamma\vdash e:\tau' \\ \tau' \sqsubseteq \tau }
	{ \Gamma\vdash_{\tau_r} x:=e \dashv\Gamma}
}

\newcommand{\TSWHILE} {
	\inferrule*[right=TS-WHILE]
	{ \Gamma \vdash e : \kw{Bool} \\ \Gamma \vdash_{\tau_r} b }
	{ \Gamma \vdash_{\tau_r} \kw{While}\ (e)\ b \dashv \Gamma }
}

\newcommand{\TSCOND} {
	\inferrule*[right=TS-COND]
	{ \Gamma\vdash e: \kw{Bool} \\ \Gamma \vdash_{\tau_r} b_1 \\ \Gamma \vdash_{\tau_r} b_2}
	{ \Gamma\vdash_{\tau_r} \kw{If}\ (e)\ b_1\ \kw{Else} \ b_2 \dashv \Gamma}
}

\newcommand{\TSRET} {
	\inferrule*[right=TS-RET]
	{\Gamma\vdash e : \tau' \\ \tau' \sqsubseteq \tau_r}
	{\Gamma \vdash_{\tau_r}\kw{Ret}\ e\dashv\Gamma}
}

\newcommand{\TSPRINT} {
	\inferrule*[right=TS-PRINT]
	{\Gamma\vdash e:\tau }
	{\Gamma \vdash_{\tau_r}\kw{Print}\ e\dashv\Gamma}
}

\newcommand{\TBBLOCK} {
	\inferrule*[right=TB-BLOCK]
	{ \Gamma_0 \vdash_{\tau_r} s_1 \dashv \Gamma_1 \\
		\Gamma_1 \vdash_{\tau_r} s_2 \dashv \Gamma_2 \\
		\cdots \\
		\Gamma_{n-1} \vdash_{\tau_r} s_n \dashv \Gamma_n }
	{ \Gamma_0 \vdash_{\tau_r} \{\, s_1; \ldots; s_n \,\} }
}

\newcommand{\TPPROG} {
	\inferrule*[right=TP-PROG]
	{ \Gamma_0 = \overline{f_j : (\tau_{j,1}, \ldots, \tau_{j,k_j}) \to \tau'_j} \\
		\big(\, \Gamma_0,\ x_{j,1}:\tau_{j,1}, \ldots, x_{j,k_j}:\tau_{j,k_j} \vdash_{\tau'_j} b_j \,\big)_{j} \\
		\Gamma_0 \vdash_{\bullet} \{\, s_1; \ldots; s_n \,\} }
	{ \vdash p\ \text{is well-typed} }
}

\newcommand{\fwidest} {
	\mathsf{widest}(\tau_1, \tau_2) =
	\begin{cases}
		\tau_1 & \tau_1 = \tau_2 \\
		\text{the wider of } \tau_1,\tau_2 & \tau_1,\tau_2 \in \{\kw{Nat},\kw{Int},\kw{Real}\} \\
		\text{undefined} & \text{otherwise}
	\end{cases}
}

\newcommand{\fneg} {
	\mathsf{neg}(\tau) =
	\begin{cases}
		\kw{Int} & \tau = \kw{Nat} \\
		\kw{Int} & \tau = \kw{Int} \\
		\kw{Real} & \tau = \kw{Real} \\
		\text{undefined} & \text{otherwise}
	\end{cases}
}

\newcommand{\fsig} {
	\mathsf{sig}(\oplus, \tau_1, \tau_2) =
	\begin{cases}
		\mathsf{widest}(\tau_1, \tau_2) & \oplus \in \{\kw{+},\kw{*},\kw{/},\kw{\%}\},\ \tau_1,\tau_2 \in \mathrm{Num} \\
		\mathsf{widest}(\mathsf{widest}(\tau_1,\tau_2),\,\kw{Int}) & \oplus = \kw{-},\ \tau_1,\tau_2 \in \mathrm{Num} \\
		\kw{Bool} & \oplus \in \{\kw{<},\kw{<=},\kw{>=},\kw{>}\},\ \tau_1,\tau_2 \in \mathrm{Num} \\
		\kw{Bool} & \oplus \in \{\kw{==},\kw{!=}\},\ \tau_1,\tau_2 \in \mathrm{Num} \text{ or } \tau_1 = \tau_2 = \kw{Bool} \\
		\kw{Bool} & \oplus \in \{\kw{\&\&},\kw{||}\},\ \tau_1 = \tau_2 = \kw{Bool} \\
		\text{undefined} & \text{otherwise}
	\end{cases}
}

\newcommand{\ENUM} {
	\inferrule*[right=E-Num]
	{ }
	{\config{n,\rho,\sigma} \eval \config{n,\sigma}}
}
\newcommand{\ETRUE} {
	\inferrule*[right=E-True]
	{ }
	{\config{\kw{True},\rho,\sigma} \eval \config{\kw{True},\sigma}}
}
\newcommand{\EFALSE} {
	\inferrule*[right=E-False]
	{ }
	{\config{\kw{False},\rho,\sigma} \eval \config{\kw{False},\sigma}}
}
\newcommand{\EVAR} {
	\inferrule*[right=E-Var]
	{ \rho(x)=\ell}
	{\config{x,\rho,\sigma} \eval \config{\sigma(\ell),\sigma}}
}
\newcommand{\ENEG} {
	\inferrule*[right=E-Neg]
	{\config{e, \rho, \sigma} \eval \config{v, \sigma'}}
	{\config{-e, \rho, \sigma} \eval \config{\mathsf{\hat{-}}v, \sigma'}}
}
\newcommand{\EOP} {
	\inferrule*[right=E-Op]
	{\config{e_1, \rho, \sigma} \eval \config{v_1, \sigma_1} \\
		\config{e_2, \rho, \sigma_1} \eval \config{v_2, \sigma_2}}
	{\config{e_1 \oplus e_2, \rho, \sigma} \eval \config{v_1 \mathbin{\hat{\oplus}} v_2, \sigma_2}}
}
\newcommand{\ECONDT} {
	\inferrule*[right=E-Cond-T]
	{\config{e,\rho, \sigma}\eval\config{\kw{True}, \sigma'} \\ \config{e_1, \rho, \sigma'}\eval\config{v_1, \sigma''}}
	{\config{e\ ?\ e_1 : e_2, \rho, \sigma}
		\eval
		\config{v_1, \sigma''}}
}
\newcommand{\ECONDF} {
	\inferrule*[right=E-Cond-F]
	{\config{e,\rho, \sigma}\eval\config{\kw{False}, \sigma'} \\ \config{e_2, \rho, \sigma'}\eval\config{v_2, \sigma''}}
	{\config{e\ ?\ e_1 : e_2, \rho, \sigma}
		\eval
		\config{v_2, \sigma''}}
}	
\newcommand{\ELAMBDA} {
	\inferrule*[right=E-Lambda]
	{ }
	{\config{\kw{Lambda}\ (\overline{x : \tau}) : \tau' \Rightarrow e, \rho, \sigma}
		\eval
		\config{\clo{\overline{x}}{\{\kw{Ret}\ e;\}}{\rho}, \sigma}}
}
\newcommand{\EAPP} {
	\inferrule*[right=E-App]
	{\config{e, \rho, \sigma} \eval \config{\clo{\overline{x}}{b}{\rho_\theta}, \sigma_0} \\
		\config{e_i, \rho, \sigma_{i-1}} \eval \config{v_i, \sigma_i} \quad (1 \le i \le n) \\\\
		\overline{\ell} \notin \mathrm{dom}(\sigma_n) \\
		\config{b, \rho_\theta[\overline{x} \mapsto \overline{\ell}], \sigma_n[\overline{\ell} \mapsto \overline{v}]} \eval \config{\mathsf{ret}\ v, \sigma'}}
	{\config{e(\overline{e}), \rho, \sigma} \eval \config{v, \sigma'}}
}
\newcommand{\SVAR} {
	\inferrule*[right=S-Var]
	{\config{e, \rho, \sigma} \eval \config{v, \sigma'} \\
		\ell \notin \mathrm{dom}(\sigma')}
	{\config{\kw{Var}\ x : \tau = e, \rho, \sigma}
		\eval \config{\mathsf{normal}, \rho[x \mapsto \ell], \sigma'[\ell \mapsto v]}}
}
\newcommand{\SREAD} {
	\inferrule*[right=S-Read]
	{v \text{ of type } \tau \text{ is read from input} \\ \ell \notin \mathrm{dom}(\sigma)}
	{\config{\kw{Var}\ x : \tau = \kw{Read}, \rho, \sigma}
		\eval \config{\mathsf{normal}, \rho[x \mapsto \ell], \sigma[\ell \mapsto v]}}
}
\newcommand{\SASSIGN} {
	\inferrule*[right=S-Assign]
	{\config{e, \rho, \sigma} \eval \config{v, \sigma'} \\
		\rho(x) = \ell}
	{\config{x := e, \rho, \sigma}
		\eval \config{\mathsf{normal}, \rho, \sigma'[\ell \mapsto v]}}
}
\newcommand{\SWHILET} {
	\inferrule*[right=S-While-T]
	{\config{e, \rho, \sigma} \eval \config{\kw{True}, \sigma'} \\
		\config{b, \rho, \sigma'} \eval \config{\mathsf{normal}, \sigma''} \\
		\config{\kw{While}\ (e)\ b, \rho, \sigma''} \eval \config{r, \rho, \sigma'''}}
	{\config{\kw{While}\ (e)\ b, \rho, \sigma} \eval \config{r, \rho, \sigma'''}}
}
\newcommand{\SWHILEF} {
	\inferrule*[right=S-While-F]
	{\config{e, \rho, \sigma} \eval \config{\kw{False}, \sigma'}}
	{\config{\kw{While}\ (e)\ b, \rho, \sigma} \eval \config{\mathsf{normal}, \rho, \sigma'}}
}
\newcommand{\SWHILERET} {
	\inferrule*[right=S-While-Ret]
	{\config{e, \rho, \sigma} \eval \config{\kw{True}, \sigma'} \\
		\config{b, \rho, \sigma'} \eval \config{\mathsf{ret}\ v, \sigma''}}
	{\config{\kw{While}\ (e)\ b, \rho, \sigma} \eval \config{\mathsf{ret}\ v, \rho, \sigma''}}
}
\newcommand{\SCONDT} {
	\inferrule*[right=S-COND-T]
	{\config{e,\rho,\sigma}\eval\config{\kw{True}, \sigma'} \\ \config{b_1, \rho,\sigma'}\eval\config{r_1,\sigma''}}
	{\config{\kw{If}\ (e)\ b_1\ \kw{Else}\ b_2, \rho, \sigma} \eval \config{r_1,\rho,\sigma''}}
}
\newcommand{\SCONDF} {
	\inferrule*[right=S-COND-F]
	{\config{e,\rho,\sigma}\eval\config{\kw{False}, \sigma'} \\ \config{b_2, \rho,\sigma'}\eval\config{r_2,\sigma''}}
	{\config{\kw{If}\ (e)\ b_1\ \kw{Else}\ b_2, \rho, \sigma} \eval \config{r_2,\rho,\sigma''}}
}
\newcommand{\SRET} {
	\inferrule*[right=S-Ret]
	{\config{e,\rho,\sigma}\eval\config{v,\sigma'}}
	{\config{\kw{Ret}\ e, \rho, \sigma}\eval \config{\mathsf{ret}\ v, \rho,\sigma'}}
}
\newcommand{\SPRINT} {
	\inferrule*[right=S-Print]
	{\config{e,\rho,\sigma}\eval\config{v,\sigma'}}
	{\config{\kw{Print}\ e, \rho, \sigma}\eval\config{\mathsf{normal}, \rho, \sigma'}}
}
\newcommand{\BBLOCK} {
	\inferrule*[right=B-Empty]
	{ }
	{\config{\{ \ \}, \rho, \sigma} \eval \config{\mathsf{normal}, \sigma}}
}
\newcommand{\BSEQ} {
	\inferrule*[right=B-Seq]
	{\config{s_1, \rho, \sigma}\eval\config{\mathsf{normal}, \rho', \sigma'} \\ \config{\{\ \overline{s}\ \}, \rho', \sigma'} \eval \config{r, \sigma''}}
	{\config{\{ s_1; \overline{s}; \}, \rho, \sigma} \eval \config{r, \sigma''}}
}
\newcommand{\BRET} {
	\inferrule*[right=B-Ret]
	{\config{s_1, \rho, \sigma}\eval\config{\mathsf{ret}\ v, \rho', \sigma'}}
	{\config{\{ s_1; \overline{s}; \}, \rho, \sigma} \eval \config{\mathsf{ret}\ v, \sigma'}}
}
\newcommand{\PPROG} {
	\inferrule*[right=P-Prog]
	{\mathrm{funcs}(p) = \kw{Func}\ F_j (\overline{x_j : \tau_j}) : \tau'_j = b_j \quad (1 \le j \le k) \\
		\mathrm{stmts}(p) = \overline{s} \\\\
		\ell_1, \ldots, \ell_k \text{ distinct} \\
		\rho_g = [F_1 \mapsto \ell_1, \ldots, F_k \mapsto \ell_k] \\\\
		\sigma_0 = [\ell_1 \mapsto \clo{\overline{x_1}}{b_1}{\rho_g}, \ldots, \ell_k \mapsto \clo{\overline{x_k}}{b_k}{\rho_g}] \\\\
		\config{\{ \overline{s} \}, \rho_g, \sigma_0} \eval \config{\mathsf{normal}, \sigma'}}
	{\config{p} \eval \config{\sigma'}}
}

\title{\langname}
\author{David Tilman \\ Student ID: \studentid}
\date{\today}

\begin{document}
\maketitle

\begin{center}
  \begin{tabular}{rl}
    Repository: & \url{\giturl} \\
    Commit:     & \texttt{\githash}
  \end{tabular}
\end{center}

% ==========================================================================
% MAIN TEXT: 3 PAGES MAX (non-extended). Put the KEY / distinctive rules here
% with prose that shows understanding; push the FULL rule sets to the appendix.
% SCOPE: IMP core + functions/lambdas/closures + numeric-subtyping type system.
% NO lists, NO records. (Borrow from the numa repo, deleting those productions.)
% ==========================================================================

% ==========================================================================
\section{Informal Description}
% 1--2 paragraphs. What the language is and how it MEANINGFULLY extends IMP
% (and borrows from the lambda-calculus) beyond the lectures. Name the two
% extensions explicitly: (1) first-class functions with closures and static
% scoping; (2) a static type system with numeric subtyping Nat <: Int <: Real.


% ==========================================================================
\section{Formal Syntax and Semantics}

\subsection{Abstract Syntax}
% BNF for: operators, types (base + function types), expressions, blocks,
% statements, declarations (Func), programs. Borrow from numa; drop the list
% and record productions.
%
\[
\begin{array}{rcll}
	  \oplus & ::=  & \kw{+} \mid \kw{-} \mid \kw{*} \mid \kw{/} \mid \kw{\%}
					   \mid \kw{<} \mid \kw{<=} \mid \kw{==} \mid \kw{!=} \mid \kw{>=} \mid \kw{>} 
					   \mid \kw{||} \mid \kw{\&\&}                                                                                     & \text{binary operators}                 \\
	\tau   & ::=  & \kw{Nat} \mid \kw{Int} \mid \kw{Real} \mid \kw{Bool} \mid  (\tau_1, \dots, \tau_n) \to \tau                        & \text{types}                            \\
	e      & ::=  & n \mid \kw{True} \mid \kw{False} \mid x 
	                  \mid \kw{Lambda}\ (x_1 : \tau_1, \dots, x_n : \tau_n) : \tau \Rightarrow e                                       & \text{literals, variable, lambda}       \\
	& \mid & -e \mid e_1 \oplus e_2 \mid e(e_1, \dots, e_n) \mid e\ ?\ e_1 : e_2                                                       & \text{operations, application, ternary} \\
	s      & ::=  & \kw{Var}\ x : \tau = e \mid \kw{Var}\ x : \tau = \kw{Read} \mid  x := e                                            & \text{declaration, read, assign}        \\
	& \mid & \kw{While}\ (e)\ b \mid \kw{If}\ (e)\ b_1\ \kw{Else}\ b_2                                                                 & \text{control flow}                     \\
	& \mid & \kw{Ret}\ e \mid \kw{Print}\ e                                                                                            & \text{return or print}                  \\
	d & ::= & \kw{Func}\ F\ (x_1 : \tau_1, \dots, x_n : \tau_n) : \tau = b                                                             & \text{function declaration}             \\
	b      & ::=  & \{\, s_1;\ \dots;\ s_n; \,\}                                                                                       & \text{block } (n \ge 0)                 \\
	c      & ::=  & d \mid s                                                                                                           & \text{top-level item}                   \\
	p      & ::=  & c_1;\ \dots;\ c_n;                                                                                                 & \text{program}
\end{array}
\]


\subsection{Type System}
% KEY rules + prose only. Lead with the interesting choices:
%   - subtyping  Nat <: Int <: Real  (and the widest/sig helpers);
%   - the context-threading statement judgement  Gamma |-_{tau_r} s -| Gamma';
%   - TE-APP / TE-LAMBDA with the subtyping premise (TE-LAMBDA extends Gamma);
%   - TP-PROG (functions mutually visible; top-level Ret untypable).
% Full typing rules -> Appendix A.
%
\noindent The type system uses three partial meta-functions, writing
$\mathrm{Num} = \{\kw{Nat}, \kw{Int}, \kw{Real}\}$ for the numeric types:
$\mathsf{neg}$ types unary negation (widening a $\kw{Nat}$ to $\kw{Int}$),
$\mathsf{widest}$ is the least upper bound of $\sqsubseteq$ (used by the ternary),
and $\mathsf{sig}$ gives the result type of a binary operator.
\[ \fneg \qquad \fwidest \]
\[ \fsig \]

\begin{mathpar}
	\STREF \and \STTRANS \and \STNATINT \and \STINTREAL \\
	\TENAT \and \TEREAL \\
	\TETRUE \and \TEFALSE \\
	\TEVAR \and \TELAMBDA \\
	\TENEGATE \and \TEOP \and \TEAPP \and \TECOND \\
	\TSDECL \and \TSREAD \and \TSASSIGN \\
	\TSWHILE \and \TSCOND \\
	\TSRET \and \TSPRINT \\
	\TBBLOCK \and \TPPROG
\end{mathpar}


\subsection{Operational Semantics}
% Notation (rho : Var -> Loc, sigma : Loc -> Value, closures, outcomes
% normal | ret v), the judgement forms, then the DISTINCTIVE rules:
%   - E-APP (closures, fresh locations, static scoping);
%   - P-PROG (hoisting: every closure captures rho_g);
%   - S-WHILE-T / S-WHILE-RET (outcome threading).
% Trivial rules (E-NUM, E-VAR, S-ASSIGN, B-SEQ, ...) -> Appendix B.
%
\begin{mathpar}
	\ENUM \and \ETRUE \and \EFALSE \and \ELAMBDA \\
	\ENEG \and \EOP \\
	\EAPP \\ 
	\ECONDT \and \ECONDF \\
	\SVAR \and \SREAD \and \SASSIGN \\
	\SWHILET \and \SWHILEF \and \SWHILERET \\
	\SCONDT \and \SCONDF \\
	\SRET \and \SPRINT \\
	\BBLOCK \and \BSEQ \and \BRET \\
	\PPROG
\end{mathpar}


% ==========================================================================
\section{Implementation}
% Technical details + match with the formal semantics. Keep tight:
%   - host language (C#); pipeline Lexer -> Parser -> Type Checker -> AST interp;
%   - sealed records mirror the grammar one-to-one;
%   - Env/Ref model rho/sigma;
%   - ONE correspondence example (e.g. E_Var) tying code to a rule;
%   - note limitations: monomorphic lambdas; numeric precision/overflow.


% ==========================================================================
\section{Distinctive-Feature Example Program}
% 1--2 paragraphs. A higher-order program -- e.g. a function that RETURNS a
% closure capturing a local variable (make_adder style). Show the source,
% state what it demonstrates, and how it uses closures + static scoping.


% ==========================================================================
\section{Reflection on AI and Collaboration}
% 1 paragraph, in your own words. Condense AI_Note.txt honestly; note any
% collaboration with other students (code + report written independently).


% ==========================================================================
\appendix

\section{Full Type System Rules}
% The complete typing rules not shown in Section 2.2.


\section{Full Operational Semantics Rules}
% The complete reduction rules not shown in Section 2.3.


% --------------------------------------------------------------------------
% REQUIRED example programs live in /example_programs (source in the repo):
%   addition, half (floor n/2), factorial, fibonacci, + the distinctive one.
% --------------------------------------------------------------------------

\end{document}
